\section{The scalar finite part and the regular model}
\label{scalar:appendix}\label{app:scalar-estimates}

The scalar finite part controls the contribution left after the central
root integral has been separated from the completed model.  We give its
periodic maximum argument first, including the derivative jump at the
end of a period.  We then retain the correlations between the endpoint
weights when applying that maximum.  All scalar assertions up to
Proposition~\ref{scalar:origin-bounds} are analytic.  The explicit
enclosure used in that proposition is a separate computer-assisted
input.

\subsection{The periodic profile and its derivative measure}

For $0<t<1$ put
\begin{equation}\label{scalar:C-definition}
 C(t)=2\gamma+\frac1t+\frac1{1-t}
 +\sum_{k=1}^{\infty}
 \left\{\frac1{k+t}+\frac1{k+1-t}-\frac2k\right\},
 \qquad C_0=C(1/2)=2\gamma+4\log2.
\end{equation}
The series converges absolutely and is symmetric under $t\mapsto1-t$.
Define the continuous periodic function
\begin{equation}\label{scalar:f-definition}
 \begin{gathered}
 \Phi(t)=\frac{\sin(\pi t)}\pi\{C(t)-2\log\pi\}
 \quad(0<t<1),\qquad \Phi(0)=\Phi(1)=1,\\
 f(x)=\Phi\left(\frac12+\frac{x}{2\pi}\right).
 \end{gathered}
\end{equation}
Here $\Phi$ is extended with period one before $f$ is evaluated.  Thus
$f$ is even and has period $2\pi$.  On the principal interval, the
symmetry of the reciprocal pairs gives
\begin{equation}\label{scalar:f-series}
 f(x)=\frac{\cos(x/2)}\pi
 \left\{k_0+4x^2\sum_{\substack{\ell\geq1\\\ell\ {\rm odd}}}
       \frac1{\ell(\ell^2\pi^2-x^2)}\right\},
 \qquad k_0=C_0-2\log\pi,\qquad |x|<\pi.
\end{equation}
These equalities also fix the normalization of $f$; in particular
$f(0)=k_0/\pi$.

\begin{lemma}[A global one-sided curvature bound]
\label{scalar:curvature}
The function $f$ is Lipschitz and its derivative has bounded variation
on the circle.  If $g=f''$ denotes its ordinary second derivative on
$(-\pi,\pi)$, then its distributional second derivative is
\begin{equation}\label{scalar:derivative-measure}
 df'=g(x)\,dx-J\delta_{\pi},\qquad
 J=\frac{2(\log\pi-\gamma)}\pi>\frac1\pi,
 \qquad df'\leq\frac3{\pi^2}\,dx.
\end{equation}
The atom is counted once on the circle.  In particular, for every real
$x,v$,
\begin{equation}\label{scalar:second-difference}
 f(x+v)+f(x-v)-2f(x)\leq\frac{3v^2}{\pi^2}.
\end{equation}
\end{lemma}

\begin{proof}
For $0<t\leq1/2$, rearrangement of
\eqref{scalar:C-definition} gives
\[
 C(t)=\frac1t+2\gamma+R(t),\qquad
 R(t)=2t^2\sum_{k\geq1}\frac1{k(k^2-t^2)}
     =2\sum_{j\geq1}\zeta(2j+1)t^{2j}.
\]
All differentiated series below converge uniformly on this half
interval.  The elementary estimate $\zeta(2j+1)\leq\zeta(3)<5/4$
follows by bounding the tail after $k=2$ by its integral.  Consequently
\[
 R'(t)\leq\frac{5t}{(1-t^2)^2},\qquad
 R''(t)\leq R''(1/2)=28\zeta(3)-16<19.
\]
The last identity also follows by differentiating the reciprocal
pairs and separating even and odd denominators.  Write
$k_{\rm end}=2\gamma-2\log\pi$.  Then
\[
 \Phi(t)=\frac{\sin(\pi t)}{\pi t}
       +k_{\rm end}\frac{\sin(\pi t)}\pi
       +R(t)\frac{\sin(\pi t)}\pi.
\]
The integral representation of the first term gives
\[
 \left(\frac{\sin(\pi t)}{\pi t}\right)''
 =-\pi^2\int_0^1u^2\cos(\pi tu)\,du
 \leq-\pi^2\left(\frac13-\frac{\pi^2}{40}\right)
 \leq-\frac{\pi^2}{12}.
\]
We use the elementary bounds
$23/40<\gamma<3/5$, $\log\pi<6/5$, and $\pi^2<10$.
They imply $-k_{\rm end}<5/4$.  Moreover,
$\cos(\pi t)\leq\pi(1/2-t)$ and
\[
 \frac{t(1/2-t)}{(1-t^2)^2}\leq\frac3{40}.
\]
Indeed, after multiplication by the positive denominator, the last
inequality is
$3+34t^2+3t^4-20t\geq0$; already the quadratic part has minimum
$3-100/34>0$.  Differentiating the three terms of $\Phi$, and dropping
only the nonpositive term $-\pi R(t)\sin(\pi t)$, now yields
\[
 \Phi''(t)\leq-\frac{\pi^2}{12}+2\pi+\frac{19}\pi<12.
\]
For the final comparison, $31/10<\pi<22/7$ suffices.
Reflection gives the same bound on the other half of an open period.
The derivative matches at the half integer by symmetry.  At an integer
its jump is $2k_{\rm end}<0$, because the displayed expansion is
$\Phi(t)=1+k_{\rm end}t+O(t^2)$ as $t\downarrow0$.
Thus $d\Phi'\leq12\,dt$, including every negative integer atom.
Scaling by $t=1/2+x/(2\pi)$ proves the measure bound and the value of
$J$ in \eqref{scalar:derivative-measure}.  The strict lower bound on
$J$ uses $\log\pi>11/10$ and $\gamma<3/5$.

The same expansions and the uniformly convergent derivatives show
that the ordinary derivatives have finite one-sided limits.  Hence
$f$ is globally Lipschitz and $f'$ has bounded variation.  Integrating
the measure bound against the nonnegative triangular kernel
$(|v|-|t|)_+$ proves \eqref{scalar:second-difference}.
\end{proof}

\begin{lemma}[Range, smooth bounds, and periodic variation]
\label{scalar:global-profile-bounds}
The periodic profile satisfies
\begin{equation}\label{scalar:profile-bounds}
 \frac12<f(x)\leq1,\qquad \operatorname{Lip}(f)<2,\qquad
 |f''(x)|<1\quad(-\pi<x<\pi),\qquad
 \operatorname{Var}_{\rm per}(f')<4.
\end{equation}
The variation includes the derivative jump at the periodic endpoint.
\end{lemma}

\begin{proof}
In the half-period representation in Lemma~\ref{scalar:curvature},
$0\leq R(t)\leq4t^2\leq1$.  Since
$2\gamma-2\log\pi+R(t)<6/5-22/10+1=0$, that representation gives
$\Phi(t)\leq\sin(\pi t)/(\pi t)\leq1$.  Reflection and continuous
periodicity give the upper bound everywhere.

For the lower bound write $0\leq s<1$, $c_s=\cos(\pi s/2)$ in
\eqref{scalar:f-series}.  Keeping its first reciprocal pair gives
\[
 \pi f(\pi s)\geq c_sk_0+\frac{4c_ss^2}{1-s^2}.
\]
The inequality $c_s\geq(3/4)(1-s^2)$ follows by putting $r=1-s$:
the function $\sin(\pi r/2)-(3/4)r(2-r)$ vanishes at zero,
has positive derivative there, and its second derivative decreases
through zero once; its first derivative consequently increases and
then decreases to zero at $r=1$.  Also
$1-c_s\leq\pi^2s^2/8$, $k_0<2$, and $\pi^2<10$.  Thus
$4c_ss^2/(1-s^2)\geq3s^2\geq k_0(1-c_s)$, proving
$f(\pi s)\geq f(0)$.  The elementary bounds
$\gamma>23/40$, $\log2>69/100$, $\log\pi<23/20$ give
$k_0>161/100>\pi/2$; hence $f(0)>1/2$.  Endpoint continuity completes
the range bound.

For the derivative estimates, the reciprocal series and integral
comparison give on $0<t\leq1/2$
\[
 R\leq1,\qquad R'\leq\frac{32}{3}t,\qquad
 R''\leq\frac{608}{27}.
\]
For example, differentiate
$R=2t^2\sum_{k\geq1}[k(k^2-t^2)]^{-1}$, use
$k^2-t^2\geq3k^2/4$, and bound the reciprocal cubes and fifth powers
by $3/2$ and $5/4$.  The derivative of the sinc term is bounded by
$\pi^2t/3<2$, the next term by $5$, and the derivative of the last
term by $R+tR'\leq11/3<4$.  Thus $|\Phi'|<12$, so
$\operatorname{Lip}(f)<6/\pi<2$.  The five second-derivative
contributions, in their displayed order and using $3<\pi<10/3$,
are bounded in absolute value by
\[
 \frac{10}{3},\quad10,\quad\frac{608}{81},\quad
 \frac{32}{3},\quad\frac{10}{3}.
\]
Their sum is $2822/81<35$.  Reflection proves this on each open
period, and scaling gives $|f''|<35/(4\pi^2)<1$.
Finally the measure $\mu=df'$ has total mass zero.  By
\eqref{scalar:derivative-measure}, its positive mass is at most
$6/\pi$, and its negative mass is equal to its positive mass.
Therefore $|\mu|(\mathbb R/(2\pi\mathbb Z))\leq12/\pi<4$.
\end{proof}

\subsection{An endpoint-cancelled integral and the scalar maximum}

Define the scalar finite part by
\begin{equation}\label{scalar:H-definition}
 H(x)=\frac1\pi\lim_{\delta\downarrow0}
 \left\{\int_{\delta}^{2\pi-\delta}
       \frac{f(x+v)}{2\sin(v/2)}\,dv
              +2f(x)\log\frac\delta\pi\right\}.
\end{equation}
Pairing the two endpoints and using
$\int_\delta^{2\pi-\delta}(2\sin(v/2))^{-1}dv
=2\log\cot(\delta/4)$ gives the absolutely convergent expression
\begin{equation}\label{scalar:H-paired}
 H(x)=\lambda_0f(x)+\frac1\pi\int_0^\pi
   \frac{f(x+v)+f(x-v)-2f(x)}{2\sin(v/2)}\,dv,
 \qquad \lambda_0=\frac2\pi\log\frac4\pi.
\end{equation}
Absolute convergence follows from the Lipschitz bound: the numerator
is $O(v)$.  Also
\begin{equation}\label{scalar:lambda-bound}
 0<\lambda_0<\frac1{2\pi}.
\end{equation}
For the upper bound use $\pi>25/8$ and
$\exp(1/4)>32/25$, the latter already following from the first four
positive terms of its exponential series.

\begin{lemma}[Strict decrease of the smooth second derivative]
\label{scalar:third-derivative}
On $0<x<\pi$,
\begin{equation}\label{scalar:f-third}
 f'''(x)<-\frac{311}{71680}\sin(x/2)<0.
\end{equation}
In particular the ordinary second derivative $g$ in
\eqref{scalar:derivative-measure} is even and decreases on $[0,\pi]$,
using its finite one-sided values at the endpoints.
\end{lemma}

\begin{proof}
Put $f_0=f(0)$.  We first remove the apparent endpoint singularity in
\eqref{scalar:f-series}.  For $\ell=2j+1$, product-to-sum integration
gives
\[
 I_\ell(x)=\int_0^{1/2}
  \{\cos(xt)-\cos(x/2)\}\cos(\ell\pi t)\,dt
 =\frac{(-1)^j\cos(x/2)x^2}
        {\ell\pi(\ell^2\pi^2-x^2)}.
\]
Thus the part of \eqref{scalar:f-series} beyond $f_0\cos(x/2)$ is
$4\sum_{j\geq0}(-1)^jI_{2j+1}(x)$, an absolutely convergent series
for $|x|<\pi$.  To exchange summation and integration, first insert
$0<r<1$.  The resulting kernel is
\[
 K_r(t)=\sum_{j\geq0}(-1)^jr^{2j+1}\cos((2j+1)\pi t)
 =\frac{r(1+r^2)\cos(\pi t)}
 {(1-r^2)^2+4r^2\cos^2(\pi t)}.
\]
For $1/2\leq r<1$ and $0\leq t<1/2$,
\[
 0\leq K_r(t)\leq\frac5{8\cos(\pi t)},\qquad
 K_r(t)\longrightarrow\frac1{2\cos(\pi t)}.
\]
The bound follows by discarding $(1-r^2)^2$ in the denominator and
using $(r+r^{-1})/4\leq5/8$.  On $|x|\leq\pi$,
\[
 |\cos(xt)-\cos(x/2)|\leq |x|(1/2-t),\qquad
 \cos(\pi t)\geq2(1/2-t).
\]
Dominated convergence, and then absolute convergence of the series of
$I_\ell$, give the exact endpoint-cancelled formula
\begin{equation}\label{scalar:f-integral}
 f(x)=f_0\cos(x/2)+2\int_0^{1/2}
          \frac{\cos(xt)-\cos(x/2)}{\cos(\pi t)}\,dt,
 \qquad |x|\leq\pi.
\end{equation}
At $x=\pi$ the integral is $1$, the prescribed continuous endpoint
value.  Evenness gives the other endpoint.

For each of the first three $x$ derivatives, the differentiated
numerator still vanishes at $t=1/2$.  Its $t$ derivative is uniformly
bounded on the compact rectangle under consideration.  The same bound
by a constant times $1/2-t$ therefore justifies differentiation under
the integral, including finite one-sided limits at $x=\pm\pi$.
In particular
\begin{equation}\label{scalar:f-third-integral}
 f'''(x)=\frac{f_0}{8}\sin(x/2)
 +2\int_0^{1/2}
  \frac{t^3\sin(tx)-(1/8)\sin(x/2)}{\cos(\pi t)}\,dt.
\end{equation}
For $0<t<1/2$, the quotient
$\sin(tx)/\sin(x/2)$ increases for $0<x\leq\pi$.
Indeed $z\cot z$ strictly decreases on $(0,\pi/2]$, since
its derivative is $(\sin z\cos z-z)/\sin^2z<0$.
Consequently
\[
 \frac{\sin(tx)}{\sin(x/2)}\leq\sin(\pi t).
\]
Divide \eqref{scalar:f-third-integral} by $\sin(x/2)>0$.  We obtain
\[
 \frac{f'''(x)}{\sin(x/2)}\leq
 B_3:=\frac{f_0}{8}
 +2\int_0^{1/2}\frac{t^3\sin(\pi t)-1/8}{\cos(\pi t)}\,dt.
\]
The numerator still vanishes at the upper endpoint.  Elementary
constants give
$f_0=(C_0-2\log\pi)/\pi<(9/5)/(25/8)<3/5$.
For $t<1/2$, the integrand numerator is negative, so division by
$\cos(\pi t)\leq1$ can only decrease it.  On $[0,1/4]$ use
$\sin(\pi t)\leq(22/7)t$, on $[1/4,3/8]$ use
$t^3\sin(\pi t)\leq27/512$, and discard the negative last interval.
It follows that
\[
 B_3<\frac3{40}
   +\left(\frac{44}{35\cdot1024}-\frac1{16}\right)
   -\frac{37}{2048}
   =-\frac{311}{71680}.
\]
This proves \eqref{scalar:f-third}.  The symmetry of $g$ follows from
the evenness of $f$.  This proof concerns the ordinary derivatives on
the principal interval; it does not remove the atom in
\eqref{scalar:derivative-measure}.
\end{proof}

\begin{corollary}[A smooth third-derivative bound]
\label{scalar:third-absolute}
On the whole open principal interval, $|f'''(x)|<13/20$.
\end{corollary}

\begin{proof}
In \eqref{scalar:f-third-integral}, the numerator
$N(t)=t^3\sin(tx)-\sin(x/2)/8$ vanishes at $t=1/2$, and
$|N'(t)|\leq3/4+\pi/8<23/20$ for $|x|\leq\pi$.
The chord bound $\cos(\pi t)\geq2(1/2-t)$ bounds
$|N(t)/\cos(\pi t)|$ by $23/40$.  After integration, the factor
two cancels the interval length $1/2$.  The exterior term is at most
$f(0)/8<3/40$, giving $23/40+3/40=13/20$.
\end{proof}

\begin{proposition}[The scalar finite part has its maximum at zero]
\label{scalar:H-maximum}
For $|x|\leq\pi$,
\begin{equation}\label{scalar:H-drop}
 H(x)-H(0)\leq-\frac{\pi-3}{4\pi^3}x^2.
\end{equation}
Thus $H(x)\leq H(0)$ for every real $x$, with equality exactly when
$x\in2\pi\mathbb Z$.
\end{proposition}

\begin{proof}
On the circle define the even kernel
\[
 W(t)=\int_{|t|}^{\pi}\log\cot(s/4)\,ds,
 \qquad |t|\leq\pi.
\]
It is bounded, nonnegative, and radially decreasing, with $W(\pi)=0$.
The triangular-measure formula for second differences, followed by
integration against $(2\sin(v/2))^{-1}$, changes
\eqref{scalar:H-paired} into
\begin{equation}\label{scalar:H-convolution}
 H(x)=\lambda_0f(x)+\frac1\pi\int_{-\pi}^{\pi}W(t)g(x+t)\,dt
                 -\frac J\pi W(\pi-|x|),\qquad |x|\leq\pi.
\end{equation}
To check the kernel, interchange the two integrals in
$\int_{|t|}^{\pi}(v-|t|)/(2\sin(v/2))\,dv$; the inner integral is
$\log\cot(s/4)$, giving $W(t)$.  Boundedness of $W$ justifies Fubini
against the total variation of $df'$.  The negative point mass in
\eqref{scalar:derivative-measure} supplies the last term in
\eqref{scalar:H-convolution}; $W(\pi)=0$ prevents double counting at
an endpoint representative.

By Lemma~\ref{scalar:third-derivative}, $g$ is even and radially
decreasing.  Subtract its finite minimum $g(\pi)$ to make it
nonnegative.  The convolution of the resulting function with $W$
is largest at zero: their layer-cake representations reduce this to
the intersection of two centered circular arcs and a translate of
one of them.  Every intersection has length at most that of the
smaller arc, and concentric arcs attain that length.  Adding the
constant $g(\pi)$ back does not affect the maximizing point.
Subtracting \eqref{scalar:H-convolution} at zero therefore gives
\[
 H(x)-H(0)\leq\lambda_0\{f(x)-f(0)\}
                    -\frac J\pi W(\pi-|x|).
\]
The measure bound and $f'(0)=0$ imply
$f(x)-f(0)\leq3x^2/(2\pi^2)$.  For $0\leq x\leq\pi$,
\[
 W(\pi-x)=\int_0^x\log\cot((\pi-s)/4)\,ds\geq\frac{x^2}{4},
\]
because the integrand vanishes at zero and its derivative is
$1/(2\cos(s/2))\geq1/2$.  Now
\eqref{scalar:lambda-bound} and $J>1/\pi$ give
\[
 H(x)-H(0)\leq
 \left(\frac3{4\pi^3}-\frac1{4\pi^2}\right)x^2.
\]
This is \eqref{scalar:H-drop}.  Its right side is negative away from
zero since $\pi>3$; periodicity completes the equality classification.
The closed endpoints are included by the displayed continuous formulas.
\end{proof}

\subsection{The endpoint weights and the full-diamond bound}

For nonnegative phases $x,y$ with $x+y\leq\pi$, put
\begin{equation}\label{scalar:phase-definitions}
 \begin{gathered}
 e=x+y,\quad d=x-y,\quad c_e=\cos(e/2),\quad c_d=\cos(d/2),\\
 a=c_ec_d=\frac{\cos x+\cos y}{2},\quad u=1-a,\quad
 z=(1-c_e)(1-c_d),\quad p=\sin x,\\
 b_0=C_0-\frac72,\quad B=\frac{C_0^2}{\pi^2},\quad
 d_0=\frac{C_0b_0}{\pi^2},\quad v_0=\frac{7b_0}{2\pi^2},\\
 g_0=\frac4{\pi^2}(\log^2\pi-\log^22),\quad
 c_{\rm b}=\frac13+g_0,\quad t_0=\frac{2\log2}{\pi}.
 \end{gathered}
\end{equation}
In particular $2-c_e-c_d=u+z$.  All weights in these formulas are
nonnegative.  Signed phases use $x=|\alpha|,y=|\beta|$ through the
separate phase comparison; no interchange of the $n$ and $n+1$ meshes
is made.

\begin{lemma}[A weighted reciprocal-pair bound]
\label{scalar:weighted-C}
With continuous endpoint products,
\begin{equation}\label{scalar:weighted-C-bound}
 \cos(\theta/2)C\left(\frac12+\frac\theta{2\pi}\right)
 \leq C_0-b_0\{1-\cos(\theta/2)\},\qquad |\theta|\leq\pi.
\end{equation}
For the endpoint contribution
$E_\infty=a/3+aC(1/2+d/(2\pi))C(1/2+e/(2\pi))/\pi^2$,
this implies the correlated estimate
\begin{equation}\label{scalar:endpoint-loss}
 E_\infty\leq\frac a3+B-d_0u-v_0z.
\end{equation}
\end{lemma}

\begin{proof}
Let $0\leq s<1$ and $c_s=\cos(\pi s/2)$.  The tangent partial
fraction identity
\[
 \frac\pi2\tan\frac{\pi s}{2}
 =\sum_{\substack{\ell\geq1\\\ell\ {\rm odd}}}
       \frac{2s}{\ell^2-s^2}
\]
follows by subtracting twice the cotangent partial fraction at $\pi s$
from the one at $\pi s/2$.  The resulting series converges absolutely
on compact subintervals of $[0,1)$.  Its derivative at zero gives
$\sum_{\ell\ {\rm odd}>0}\ell^{-2}=\pi^2/8$.
For $Q(s)=(1-s^2)/c_s$,
\[
 \frac{Q'(s)}{Q(s)}=
 \sum_{\substack{\ell\geq3\\\ell\ {\rm odd}}}
       \frac{2s}{\ell^2-s^2}
 \geq2s\left(\frac{\pi^2}{8}-1\right).
\]
Since $Q(0)=1$, integration and $\exp t\geq1+t$ imply
\[
 \frac{c_ss^2}{(1-c_s)(1-s^2)}\leq\frac8{\pi^2},
 \qquad 0<s<1.
\]
The reciprocal-pair series gives
\[
 C(1/2+s/2)-C_0
 =4s^2\sum_{\ell\ {\rm odd}>0}\frac1{\ell(\ell^2-s^2)},
 \qquad
 \sum_{\ell\ {\rm odd}\geq3}\frac1{\ell(\ell^2-1)}
 =\frac34-\log2.
\]
For the second identity, use
$1/(\ell(\ell^2-1))=(1/2)(1/(\ell-1)+1/(\ell+1))-1/\ell$
and take symmetric harmonic partial sums.  Hence
\[
 \frac{c_s\{C(1/2+s/2)-C_0\}}{1-c_s}
 \leq\frac{32}{\pi^2}\left(\frac74-\log2\right)<\frac72.
\]
The strict comparison follows from $\log2>69/100$ and
$\pi^2>197/20$.  Symmetry handles negative $s$, and at $s=0$ the
increment vanishes.  At $s\uparrow1$ the weighted increment tends
to $\pi<7/2$, as its $\ell=1$ term shows.  This proves
\eqref{scalar:weighted-C-bound} on the closed interval.

The two factors bounded in \eqref{scalar:weighted-C-bound} are
positive, including their endpoint limits.  Multiplying their upper
bounds gives exactly
\[
 [C_0-b_0(1-c_e)][C_0-b_0(1-c_d)]
 =C_0^2-C_0b_0u-\frac72b_0z.
\]
Division by $\pi^2$ and addition of $a/3$ proves
\eqref{scalar:endpoint-loss}.
\end{proof}

The completed model used in the finite-row comparisons is now specified
in full.  In the notation of Appendix~\ref{app:row-estimates}, write
$\mathcal B(v)=B(v)$ for the constant part of the row profile, to
distinguish it from the constant $B=C_0^2/\pi^2$ above.  Put
\begin{equation}\label{scalar:finitepart-normalization}
 \begin{gathered}
 U_0=c_ef(d)+c_df(e),\qquad
 D_0=\frac{U_0}{\pi}+\frac{4a}{\pi^2}\log\pi,\\
 K_\delta=\frac{4a}{\pi^2}\log^2\frac\delta\pi
       +2D_0\log\frac\delta\pi
       +\frac1\pi\int_\delta^{2\pi-\delta}\mathcal B(v)\,dv,\\
 E_\infty=\frac a3+
  \frac{a}{\pi^2}C\left(\frac12+\frac d{2\pi}\right)
                       C\left(\frac12+\frac e{2\pi}\right),\\
 C_\infty=E_\infty+\lim_{\delta\downarrow0}K_\delta.
 \end{gathered}
\end{equation}
Every product at an endpoint uses its continuous weighted value.  The
finite-part limit exists uniformly in the closed phase diamond.  Indeed,
with $v_L(t)=t$, $v_R(t)=2\pi-t$ and
\[
 Z_s(t)=\mathcal B(v_s(t))-
                 \frac{U_0+(4a/\pi)\log t}{t},
\]
the combined endpoint residual bound in the finite-row estimates is
\[
 |Z_s(t)|\leq7+\frac{p+q}{\pi}
                       \left(3+2\log\frac{4\pi}{t}\right).
\]
Subtracting two finite renormalizations gives exactly
\[
 K_\varepsilon-K_\delta=
       \frac1\pi\int_\varepsilon^\delta(Z_L+Z_R)(t)\,dt.
\]
The displayed integrable envelope, uniform for $p+q\leq2$, makes this
a uniformly Cauchy family.  The central logarithm has an integrable
envelope as in Lemma~\ref{lem:generic-cancellation}; hence the finite
integrals are continuous in phase, and their uniform limit is
continuous as well.

For every real $L$, define
\begin{equation}\label{scalar:completed-model}
 \begin{gathered}
 M_\infty(x,y,L)=\frac{4a}{\pi^2}L^2+b_{\rm full}(x,y)L
                         +C_\infty(x,y),\\
 b_{\rm full}=\frac{2U_0}{\pi}
  +\frac{8a\log2+4a\log a+2(x\sin x+y\sin y)}{\pi^2},
 \qquad 0\log0=0.
 \end{gathered}
\end{equation}
This fixes all coefficients in the error $S_n-M_\infty$.
To check the logarithmic coefficient against the row normalization,
put $f_s(t)=A(v_s(t))-4a/(\pi t)$.  Joining the endpoint and
interior coefficient integrals in the exact telescope gives
\[
 b_{\rm full}=2D_0+\frac1\pi\int_0^\pi(f_L+f_R)(t)\,dt.
\]
Here is its elementary evaluation.  Set $X=\max(x,y)$,
$Y=\min(x,y)$, so $\sin X\geq\sin Y$.  If $a>0$, the exact max
formulas for $A$ give, with $T_t=2a\cot(t/2)$,
\[
 \pi(A(v_L(t))+A(v_R(t)))=
 \begin{cases}
 2T_t,&0<t<\pi-X-Y,\\
 T_t+p+q,&\pi-X-Y<t<\pi-X+Y,\\
 2\sin X,&\pi-X+Y<t<\pi.
 \end{cases}
\]
Indeed $T_{\pi-X-Y}=p+q$ and
$T_{\pi-X+Y}=|p-q|$.  Integrate the three pieces, using
$\int T_t\,dt=4a\log\sin(t/2)$ and
$\cos((X+Y)/2)\cos((X-Y)/2)=a$, and subtract $8a/t$ at the
lower endpoint.  The result is
\[
 \int_0^\pi(f_L+f_R)(t)\,dt
 =\frac{8a\log(2/\pi)+4a\log a+2(xp+yq)}\pi.
\]
This is precisely \eqref{scalar:completed-model}.  For continuity at
$a=0$, the same max formulas give
\begin{align*}
 \pi f_L&=\max(q,p-T_t)+w(t),\\
 \pi f_R&=\max(-q,p-T_t,q-2T_t)+w(t),\\
 w(t)&=T_t-4a/t\in[-4a/\pi,0].
\end{align*}
They bound $|f_s|<1$ uniformly.  Dominated convergence and
$a\log a\to0$ extend the evaluation to that boundary, including the
two vertices.

\begin{lemma}[Loss in the logarithmic coefficients]
\label{scalar:coefficient-loss}
Let $k_s=\gamma+\log(8/\pi)$, $b_c=8k_s/\pi^2$, and
$C_*=C_\infty(0,0)$.
On the closed phase diamond,
\begin{equation}\label{scalar:b-loss}
 b_{\rm full}-b_c\leq\frac{3u}{\pi^2},\qquad
 \frac{4a}{\pi^2}\leq\frac4{\pi^2},\qquad
 b_{\rm full}<\frac{79}{5\pi^2}.
\end{equation}
For $L\geq0$, if $T=(4L^2-3L)/\pi^2$, then
\begin{equation}\label{scalar:model-loss}
 M_\infty(0,0,L)-M_\infty(x,y,L)
 \geq uT-\{C_\infty(x,y)-C_*\}.
\end{equation}
\end{lemma}

\begin{proof}
In the open phase triangle put
$D=C(1/2+e/(2\pi))+C(1/2+d/(2\pi))-2C_0$.
Lemma~\ref{scalar:weighted-C} gives
\[
 aD\leq\frac72\{c_d(1-c_e)+c_e(1-c_d)\}
 \leq\frac72u,
\]
where the last step uses $c_e+c_d\leq1+a$.
Substituting $U_0=a(2C_0+D-4\log\pi)/\pi$ in
\eqref{scalar:completed-model} gives the exact identity
\[
 b_{\rm full}-b_c=\frac2{\pi^2}
 \{aD+2a\log a+x\sin x+y\sin y-4k_su\}.
\]
For $0\leq t\leq\pi$, the inequality $\tan(t/2)\geq t/2$
gives $t\sin t\leq2(1-\cos t)$, so the two sine terms sum to at
most $4u$.  Also $a\log a\leq0$ and $k_s>3/2$.  For the last
bound one can use $\gamma>23/40$, $8/\pi>28/11$, and
\[
 \log\frac{28}{11}
 >2\left(\frac{17}{39}+\frac13\left(\frac{17}{39}\right)^3\right)
 >\frac{37}{40}.
\]
Thus $b_{\rm full}-b_c\leq(2/\pi^2)(7/2+4-6)u=3u/\pi^2$.
All weighted terms extend continuously to the boundary.  Finally,
$k_s<8/5$ follows from $\gamma<3/5$, $\pi>3$, and
$\exp(1)>8/3$.  Hence $b_{\rm full}\leq b_c+3/\pi^2
<(64/5+3)/\pi^2$.  Subtracting the two quadratic polynomials in
\eqref{scalar:completed-model} proves \eqref{scalar:model-loss}.
\end{proof}

\begin{lemma}[The global cancellation contribution]
\label{scalar:global-cancellation}
For the completed constant, the pole contribution is
$c_eH(d)+c_dH(e)$, and the remaining cancellation contribution
$K_{I,0}$ satisfies
\begin{equation}\label{scalar:global-KI}
 C_\infty=E_\infty+c_eH(d)+c_dH(e)+K_{I,0},\qquad
 K_{I,0}\leq-c_{\rm b}a+t_0p.
\end{equation}
These formulas hold on the entire closed phase diamond.
\end{lemma}

\begin{proof}
We keep the two endpoint counterterms together.  For $\varepsilon>0$
the cancellation part of the completed finite part is
\[
 K_{I,\varepsilon}=
 \frac{4a}{\pi^2}(\log^2\varepsilon-\log^2\pi)
       +\frac1{\pi^2}\int_\varepsilon^\pi
                        \{I_L(t)+I_R(t)\}\,dt.
\]
The other part is exactly
$c_eH_\varepsilon(d)+c_dH_\varepsilon(e)$, where $H_\varepsilon$
is the expression inside the limit in \eqref{scalar:H-definition}.
It therefore tends to the stated pole contribution.

For $0<t<\pi$ let $T=2a\cot(t/2)$ and
$\epsilon=\min(t,\pi-t)$.  The exact cancellation formulas of
Lemma~\ref{lem:generic-cancellation} give the two parameters
\[
 I_L=J(p/2,(q+T)/2,\epsilon),\quad
 I_R=J(p/2,|q-T|/2,\epsilon),\qquad q=\sin y.
\]
In those formulas
$J(v,s,\epsilon)=2(s-v)_+\log\sin\epsilon+R(v,s,\epsilon)$,
and $R\leq0$.  For $s\leq v$ this follows from its entropy formula;
for $s>v$, its nonnegative separation derivative and
$R(v,s,\pi/2)=0$ give the same conclusion.  Thus
\[
 I_L+I_R\leq\{(q+T-p)_++(|q-T|-p)_+\}\log\sin t.
\]
The coefficient in braces is at least $2\max(q,T)-2p\geq2T-2p$.
Since $\log\sin t\leq0$, multiplication reverses that inequality:
\begin{equation}\label{scalar:paired-I-upper}
 I_L(t)+I_R(t)\leq4a\cot(t/2)\log\sin t
                          +2p|\log\sin t|.
\end{equation}
There is no near-origin restriction here, and zero parameters use
the continuous values in the complete cancellation formulas.

The identity
$\int_0^\pi|\log\sin t|\,dt=\pi\log2$ follows by reflection,
the double-angle formula, and the substitution $t=2s$.
An antiderivative of the first term on the right of
\eqref{scalar:paired-I-upper} is
\[
 \begin{gathered}
 4a\log^2(2\sin(t/2))-4a\log^22
                    -2a\operatorname{Li}_2(\sin^2(t/2)),\\
 \operatorname{Li}_2(w)=\sum_{k\geq1}\frac{w^k}{k^2}.
 \end{gathered}
\]
Its value at $t=\pi$ is $-a\pi^2/3$, since
$\operatorname{Li}_2(1)=\pi^2/6$.
At the lower endpoint, the displayed logarithmic counterterm leaves
exactly $-c_{\rm b}a$ as $\varepsilon\downarrow0$:
$\log(2\sin(\varepsilon/2)/\varepsilon)=O(\varepsilon^2)$,
its product with $\log\varepsilon$ tends to zero, and the dilogarithm
tends to zero.  The second term contributes $t_0p$.  This proves
\eqref{scalar:global-KI}.  The original weighted finite-part integrals
and their subtracted endpoint envelopes are continuous; taking their
limits proves the statement also on the axes and at $a=0$.
\end{proof}

\subsection{The origin constant and its explicit arithmetic input}

Set $H_0=H(0)$, $j=H_0+d_0-g_0$, and $C_*=C_\infty(0,0)$.
At the origin the central root correction vanishes, and the exact
central decomposition \eqref{eq:exact-central-decomposition} gives
\begin{equation}\label{scalar:origin-identity}
 \begin{gathered}
 C_*=B-g_0+2H_0
 =\frac4{\pi^2}\left(\gamma+\log\frac8\pi\right)^2+\kappa,\\
 H_0=\frac\kappa2+\frac{2k_0}{\pi^2}\log\frac2\pi.
 \end{gathered}
\end{equation}
Here is the identification with the actual corner, including its
normalization.  At zero, $a=1$, $U_0=2k_0/\pi$, and
\eqref{scalar:completed-model} gives $b_{\rm full}=b_c$.
Lemma~\ref{lem:phase-row-global} has a right side tending to zero as
$n\to\infty$; at this phase $S_n(0,0)=\Lambda_c(n)$.  Consequently
\[
 C_*=\lim_{n\to\infty}
 \left\{\Lambda_c(n)-\frac4{\pi^2}\log^2n-b_c\log n\right\}.
\]
Theorem~\ref{thm:corner-expansion} evaluates this limit as the second
expression in \eqref{scalar:origin-identity}, since changing
$\log(n+1/2)$ to $\log n$ changes no constant in the limit.  Inserting
$2k_s=C_0+2\log2-2\log\pi$ in the first equality then gives the
stated formula for $H_0$.  These are identities of limiting constants;
they do not assert a finite lower bound obtained by truncating the
corner asymptotic expansion.

\begin{proposition}[The numerical input and its rational consequences]
\label{scalar:origin-bounds}
The scalar consequences used below require the following inherited
computer-assisted enclosure of the constants in the corner expansion:
\begin{align}
 0.919710118388492535099518225790
 &<\kappa<0.919710118388492535099518225791,
                                                     \label{scalar:kappa-enclosure}\\
 1.846160310423926079759838737947
 &<C_*<1.846160310423926079759838737948.
                                                     \label{scalar:Cstar-enclosure}
\end{align}
With this input,
\begin{equation}\label{scalar:rational-bounds}
 \begin{gathered}
 \frac7{25}<H_0<\frac{1633}{5125}<\frac8{25},\qquad
 d_0>\frac{39}{250},\quad g_0<\frac{43}{125},\\
 j>\frac{23}{250},\quad 0<t_0<\frac{221}{500}.
 \end{gathered}
\end{equation}
The enclosures \eqref{scalar:kappa-enclosure}--\eqref{scalar:Cstar-enclosure}
are the explicit computational premise, historically recorded as
PD-T271.  Their trusted base consists of outward Arb arithmetic, the
analytic tail bound, and the independent verification recorded with
that enclosure.  They are not consequences of the analytic scalar
maximum proposition.
\end{proposition}

\begin{proof}
We spell out the rational reductions so that the precision required
from the enclosure is visible.  The elementary constants are
\begin{gather*}
 \frac{161}{100}<k_0<\frac95,\qquad
 \frac{11}{10}<\log\pi<\frac{23}{20},\\
 \frac{69}{100}<\log2<\frac{6932}{10000},\qquad
 \frac{157}{50}<\pi<\frac{22}{7},\qquad
 \frac{197}{20}<\pi^2<10.
\end{gather*}
For example, $\log2>69/100$ follows from the first three terms of
$2\operatorname{atanh}(1/3)$, and its positive geometric remainder
gives the stated upper bound.  The exponential series through degree
five at $23/20$ exceeds $22/7$, proving the upper bound on
$\log\pi$.  Together with $23/40<\gamma<3/5$ these bounds give
the stated bounds on $k_0$ and $391/100<C_0<4$.

Since $\log(\pi/2)<23/50$, the first identity in
\eqref{scalar:origin-identity} and \eqref{scalar:kappa-enclosure} give
\[
 H_0>\frac9{20}-\frac{828}{4925}
       =\frac{5553}{19700}>\frac7{25}.
\]
Conversely, $\log s>2(s-1)/(s+1)$ for $s>1$ gives
$\log(\pi/2)>\log(25/16)>18/41$.  Hence
\[
 H_0<\frac{23}{50}-\frac{161}{500}\frac{18}{41}
       =\frac{1633}{5125}<\frac8{25}.
\]
Also $C_0>39/10$, $b_0>2/5$, and $\pi^2<10$ give
$d_0>39/250$, while
\[
 g_0<\frac{80}{197}
 \left\{\left(\frac{23}{20}\right)^2-
               \left(\frac{69}{100}\right)^2\right\}
 =\frac{677120}{1970000}<\frac{43}{125}.
\]
Therefore $j>7/25+39/250-43/125=23/250$.
Finally,
$t_0<6932/15700<221/500$.  These reductions use only exact rational
comparisons once the two strict enclosures have been supplied.
\end{proof}

\begin{proposition}[Correlated global master bound]
\label{scalar:master}
On the entire closed phase diamond,
\begin{align}
 C_\infty(x,y)-C_*
 &\leq-ju+t_0p-(H_0+v_0)z
 \leq-ju+t_0p,                                      \label{scalar:global-master}\\
 C_\infty(x,y)&<\frac94.                             \label{scalar:global-nine-fourths}
\end{align}
For $L\geq0$ the completed model therefore obeys
\begin{equation}\label{scalar:outer-master}
 M_\infty(0,0,L)-M_\infty(x,y,L)
 \geq u(T+j)-t_0p,\qquad T=\frac{4L^2-3L}{\pi^2}.
\end{equation}
On the closed central domain $x+y\leq\pi/2$, write
$\operatorname{Reg}=C_\infty-E_c$ with $E_c$ as in
\eqref{eq:Ec-definition}.  Then
\begin{equation}\label{scalar:regular-loss}
 \operatorname{Reg}(x,y)-\operatorname{Reg}(0,0)
 \leq-ju-(H_0+v_0)z\leq0.
\end{equation}
\end{proposition}

\begin{proof}
Apply $H\leq H_0$ in \eqref{scalar:global-KI} and add
\eqref{scalar:endpoint-loss}.  This gives
\[
 C_\infty\leq B-d_0+(d_0-g_0)a-v_0z
                      +H_0(c_e+c_d)+t_0p.
\]
Subtract $C_*=B-g_0+2H_0$ and use
$2-c_e-c_d=u+z$.  This is exactly
\eqref{scalar:global-master}, including the mixed loss.  Its last
inequality uses $H_0+v_0>0$.
Combining it with \eqref{scalar:model-loss} proves
\eqref{scalar:outer-master}.

For completeness we verify the uniform $9/4$ bound with a rational
margin.  More generally, if $H\leq H_0+\epsilon$ for
$\epsilon\geq0$, the same computation and
$a\leq(1+\cos x)/2$ give
\[
 C_\infty\leq C_*+2\epsilon-
 \frac{j+\epsilon}{2}
       +\sqrt{\frac{(j+\epsilon)^2}{4}+t_0^2}.
\]
Put $\delta=9/4-C_*$.  This upper bound is below $9/4$ whenever
\[
 \delta-2\epsilon>0,\qquad
 (\delta-2\epsilon)^2+(\delta-2\epsilon)(j+\epsilon)>t_0^2.
\]
These conditions follow by moving the square root and squaring only
after the other side is positive.  The already proved maximum permits
the weaker bound with $\epsilon=1/400$.  By
\eqref{scalar:Cstar-enclosure}, $\delta>403/1000$; thus
$\delta-2\epsilon>199/500$ and $j+\epsilon>189/2000$.
The left side of the last test is greater than
$196015/1000000$, whereas $t_0^2<195364/1000000$.
The strict rational margin is $651/1000000$.

On the central domain the exact decomposition
\eqref{eq:exact-central-decomposition} replaces the inequality for
$K_{I,0}$ by $-c_{\rm b}a+E_c$.  Subtracting its origin value in
the same way gives \eqref{scalar:regular-loss}; the global $t_0p$
allowance is absent there because the complete central correction is
retained explicitly.  All formulas use continuous weighted values,
so their stated closed boundaries are included.
\end{proof}
